% ============================================================================
%  Problems: พหุนาม (Polynomials)
%  Ordered from easy (basic) → intermediate.
%  Shared by exercise.tex and answer-key.tex
% ============================================================================

% --- Part 1: พื้นฐาน (Basic) ---
\examsection{ตอนที่ 1: พื้นฐาน — นิยามและการดำเนินการ}

\problem{
  จงหาดีกรี สัมประสิทธิ์นำ และพจน์คงที่ของพหุนาม
  $P(x) = 5x^4 - 2x^3 + 7x^2 - 9x + 3$
}{
  ดีกรี $= 4$, สัมประสิทธิ์นำ $= 5$, พจน์คงที่ $= 3$
}

\problem{
  นิพจน์ต่อไปนี้เป็นพหุนามหรือไม่?
  \begin{itemize}
    \item[ก.] $3x^2 + \sqrt{5}x - 1$
    \item[ข.] $x^{-1} + 4x$
    \item[ค.] $\frac{x^2 + 1}{2}$
  \end{itemize}
}{
  ก. \textbf{เป็น} — สัมประสิทธิ์ $\sqrt{5}$ ไม่ได้ทำให้เลขชี้กำลังเปลี่ยน

  ข. \textbf{ไม่เป็น} — $x^{-1}$ มีเลขชี้กำลังติดลบ

  ค. \textbf{เป็น} — $\frac{1}{2}(x^2+1) = 0.5x^2 + 0.5$ เลขชี้กำลังเป็นจำนวนเต็มไม่ลบ
}

\problem{
  จงหาผลลัพธ์: $(2x^3 - 5x^2 + 4x - 1) + (x^3 + 3x^2 - 2x + 7)$
}{
  $= (2+1)x^3 + (-5+3)x^2 + (4-2)x + (-1+7)$

  $= 3x^3 - 2x^2 + 2x + 6$
}

\problem{
  จงหาผลลัพธ์: $(3x^3 + x^2 - 5) - (x^3 - 4x^2 + 2x - 3)$
}{
  $= 3x^3 + x^2 - 5 - x^3 + 4x^2 - 2x + 3$

  $= 2x^3 + 5x^2 - 2x - 2$
}

\problem{
  จงหาผลคูณ: $(x + 2)(x^2 - 3x + 5)$
}{
  $= x(x^2 - 3x + 5) + 2(x^2 - 3x + 5)$

  $= x^3 - 3x^2 + 5x + 2x^2 - 6x + 10$

  $= x^3 - x^2 - x + 10$
}

% --- Part 2: การประยุกต์ (Intermediate) ---
\examsection{ตอนที่ 2: การประยุกต์ — ทฤษฎีบทเศษเหลือ, ตัวประกอบ, ทวินาม}

\problem{
  จงหาเศษจากการหาร $P(x) = x^3 - 5x^2 + 3x + 4$ ด้วย $x - 3$
  โดยใช้ทฤษฎีบทเศษเหลือ
}{
  $P(3) = 27 - 5(9) + 3(3) + 4 = 27 - 45 + 9 + 4 = -5$

  เศษ $= -5$
}

\problem{
  กำหนด $P(x) = x^3 + kx^2 - 7x + 6$ ถ้า $(x-1)$ เป็นตัวประกอบของ $P(x)$
  จงหาค่าของ $k$
}{
  $(x-1)$ เป็นตัวประกอบ ดังนั้น $P(1) = 0$

  $1 + k - 7 + 6 = 0 \;\Rightarrow\; k = 0$
}

\problem{
  จงแยกตัวประกอบของพหุนาม $P(x) = x^3 - 3x^2 - 4x + 12$ โดยการจับกลุ่ม
}{
  $x^3 - 3x^2 - 4x + 12 = x^2(x-3) - 4(x-3)$

  $= (x-3)(x^2 - 4) = (x-3)(x-2)(x+2)$
}

\problem{
  จงใช้ทฤษฎีบทตัวประกอบเพื่อแยกตัวประกอบของ $P(x) = x^3 - 4x^2 + x + 6$
}{
  ทดสอบ: $P(-1) = -1 - 4 - 1 + 6 = 0$ $\rightarrow$ $(x+1)$ เป็นตัวประกอบ

  หารด้วย $(x+1)$: ผลหาร $= x^2 - 5x + 6$

  แยกต่อ: $x^2 - 5x + 6 = (x-2)(x-3)$

  $P(x) = (x+1)(x-2)(x-3)$
}

\problem{
  จงกระจาย $(2x - 3)^5$ โดยใช้ทฤษฎีบททวินาม (เขียนเฉพาะ 3 พจน์แรกและพจน์สุดท้าย)
}{
  $(2x-3)^5 = \binom{5}{0}(2x)^5 + \binom{5}{1}(2x)^4(-3) + \binom{5}{2}(2x)^3(-3)^2 + \cdots$

  $= 32x^5 - 240x^4 + 720x^3 - 1080x^2 + 810x - 243$
}

% --- Part 3: ระดับ สอวน. (POSN Level) ---
\examsection{ตอนที่ 3: ระดับข้อสอบ สอวน.}

\problem{
  (ข้อสอบจริงปี 68 ข้อ 7) สัมประสิทธิ์ของพจน์ $x^3y^7$ จากการกระจาย $(2x+y)^{10}$ เท่ากับเท่าไร
}{
  $960$
}

\problem{
  (ข้อสอบจริงปี 67 ข้อ 9) ถ้า $(ax+3)(x-b)=3x^2-cx+9$ แล้ว $a+b-c$ เท่ากับเท่าไร
}{
  $12$
}

\problem{
  (ข้อสอบจริงปี 66 ข้อ 26) กำหนด $P(x)$ เป็นพหุนมดีกรี 3 ที่สอดคล้องกับเงื่อนไข $P(1)=P(2)=P(3)=1$ และ $P(4)=2$ จงหา $P(5)$
}{
  $5$
}

\problem{
  (ข้อสอบจริงปี 66 ข้อ 27) จงหาเศษจากการหารพหุนาม $x^{100}+x^{99}+x^{98}+...+x^2+x+1$ หารด้วย $x^2-1$
}{
  $50x+51$
}